\parttitle{I}{Foundations of Classical Reservoir Computing}

\section{The reservoir computing paradigm}
\label{sec:rc-paradigm}

Reservoir computing (RC) is a way of using a fixed dynamical system as a
computational resource. A time-dependent input drives the system; the system's
internal state, by virtue of its own dynamics, becomes a rich nonlinear
function of the input history; and a simple \emph{linear} map from the state to
the output is the only component that is ever trained. The paradigm emerged
twice independently around 2001--2002: as \emph{echo state networks} (ESNs) in
the machine-learning community~\citep{jaeger2001echo,jaeger2004harnessing} and
as \emph{liquid state machines} in computational
neuroscience~\citep{maass2002realtime}; the two were later unified under the
reservoir computing name~\citep{verstraeten2007experimental,lukosevicius2009reservoir}.

The appeal of the paradigm is threefold. First, training is a linear
least-squares problem --- convex, closed-form, and cheap --- so the pathologies of
training recurrent networks by gradient descent (vanishing and exploding
gradients through time) never arise. Second, because the dynamical core is
fixed, \emph{any} sufficiently rich dynamical system can serve as the
reservoir; this is what opens the door to physical and, in this document,
quantum reservoirs. Third, despite the simplicity of training, reservoir
computers are provably universal approximators for an important class of
input--output maps (Section~\ref{sec:universality}) and are competitive in
practice on chaotic time-series prediction, a task family directly relevant to
weather and climate~\citep{pathak2018model,arcomano2020machine,vlachas2020backpropagation}.

Figure~\ref{fig:rc-blocks} fixes the three-layer anatomy used throughout this
document: an \emph{encoding} (input layer), a fixed \emph{reservoir}, and a
trained linear \emph{read-out}. Every architecture in Part~III, quantum or
classical, is an instance of this anatomy; design work happens almost entirely
in the first two layers, and honesty in evaluation happens in the third.

\begin{figure*}[t]
\centering
\begin{tikzpicture}[node distance=9mm and 13mm]
  \node[block] (in) {input\\$u_t$};
  \node[block, right=of in, minimum width=34mm, minimum height=16mm] (res)
    {fixed reservoir\\$x_t = F(x_{t-1}, u_t)$};
  \node[block, right=of res] (out) {read-out\\$\hat y_t = W_{\mathrm{out}}\,\tilde x_t$};
  \draw[arrow] (in) -- node[above,font=\scriptsize]{$W_{\mathrm{in}}$} (res);
  \draw[arrow] (res) -- (out);
  \draw[arrow] (res.20) to[out=55,in=125,looseness=5] node[above,font=\scriptsize]{$W_{\mathrm{res}}$} (res.160);
  \node[below=2mm of in, font=\scriptsize, align=center] {given};
  \node[below=2mm of res, font=\scriptsize, align=center] {fixed (never trained)};
  \node[below=2mm of out, font=\scriptsize, align=center] {trained (linear only)};
\end{tikzpicture}
\caption{The three-layer anatomy of reservoir computing. Only the read-out is
trained; the recurrent core is fixed at initialisation. In quantum reservoir
computing (Part~II) the middle block becomes an input-driven quantum channel
and the read-out acts on measured expectation values.}
\label{fig:rc-blocks}
\end{figure*}

\section{Echo state networks}
\label{sec:esn}

\subsection{State update and read-out}

The standard ESN maintains a state $x_t \in \R^{N_x}$ updated as
\begin{equation}
\begin{aligned}
  x_{t+1} \;=\;{}& (1-\alpha)\, x_t
  + \alpha\, f\!\left(W_{\mathrm{res}}\, x_t + W_{\mathrm{in}}\, u_{t+1} + b\right),\\
  & \alpha \in (0,1],
\end{aligned}
  \label{eq:leaky}
\end{equation}
whose leak rate $\alpha$ sets the intrinsic time scale of the reservoir and is
the single most important hyperparameter to match to the autocorrelation time
of the data (a point that returns, in quantum form, as the injection interval
$\Delta t$ in Parts~II and~VIII). The output is linear in the (possibly
augmented) state $\tilde x_t = (x_t^{\!\top}, 1)^{\!\top}$:
\begin{equation}
  \hat y_t \;=\; W_{\mathrm{out}}\, \tilde x_t .
  \label{eq:readout}
\end{equation}

Three conventions from Eq.~\eqref{eq:esn-update} onward hold throughout this
document. First, a \emph{washout}: the first $T_w$ states, which still remember
the arbitrary initial condition, are discarded before training. Second,
\emph{chronological splits}: training data strictly precede validation data,
which strictly precede test data; shuffled splits leak information through
autocorrelation and are one of the most common sources of inflated results in
time-series work. Third, all errors are reported as the normalised mean-square
error
\begin{equation}
  \NMSE \;=\; \frac{\sum_t \left(y_t - \hat y_t\right)^2}{\sum_t \left(y_t - \bar y\right)^2},
  \label{eq:nmse}
\end{equation}
for which the constant predictor $\hat y_t = \bar y$ scores exactly $1$.

\subsection{Ridge regression: derivation and practice}
\label{sec:ridge}

Collect the post-washout training states row-wise into
$X \in \R^{T\times (N_x+1)}$ and the targets into $y \in \R^{T}$. Ordinary
least squares minimises $\lVert Xw - y\rVert_2^2$, but reservoir states are
strongly collinear (neighbouring neurons see overlapping input histories), so
the normal-equations matrix $X^{\!\top}X$ is ill-conditioned and the unpenalised
solution overfits noise in the state trajectory. Ridge (Tikhonov) regression
minimises instead
\begin{equation}
  J(w) \;=\; \lVert Xw - y\rVert_2^2 + \lambda \lVert w\rVert_2^2,
  \qquad \lambda > 0 .
  \label{eq:ridge-objective}
\end{equation}
Setting $\nabla_w J = 2X^{\!\top}(Xw-y) + 2\lambda w = 0$ gives the closed form
\begin{equation}
  w^\star \;=\; \left(X^{\!\top}X + \lambda \id\right)^{-1} X^{\!\top} y .
  \label{eq:ridge-solution}
\end{equation}
Since $X^{\!\top}X \succeq 0$, the regularised matrix
$X^{\!\top}X + \lambda\id \succ 0$ is always invertible, and the objective
\eqref{eq:ridge-objective} is strictly convex, so $w^\star$ is the unique global
minimiser. In the singular-value decomposition $X = U\Sigma V^{\!\top}$ with
singular values $s_i$, Eq.~\eqref{eq:ridge-solution} becomes
\begin{equation}
  w^\star \;=\; \sum_i \frac{s_i}{s_i^2 + \lambda}\, \left(u_i^{\!\top} y\right) v_i ,
  \label{eq:ridge-svd}
\end{equation}
which exhibits ridge as a \emph{spectral filter}: directions of the state
space with $s_i^2 \gg \lambda$ pass through essentially untouched, while
directions with $s_i^2 \ll \lambda$ --- precisely the poorly excited, noise-dominated
ones --- are shrunk toward zero. This picture matters doubly in the quantum
setting, where each feature carries shot noise of known magnitude
(Part~IV): $\lambda$ is then not a nuisance parameter but the dial that trades
the reservoir's expressivity against its measurement noise. The penalty is
selected on held-out data --- either a chronological validation block or, when
that block is too short to be representative, by generalised cross-validation
on the training span; it is never tuned on the test set.

\begin{algorithm*}[t]
\caption{Echo state network training (batch, ridge read-out)}
\label{alg:esn}
\begin{algorithmic}[1]
\State draw $W_{\mathrm{in}}, W_{\mathrm{res}}, b$ at random; rescale $W_{\mathrm{res}}$ to spectral radius $\rho_{\mathrm{sr}}<1$
\State $x_0 \gets 0$
\For{$t = 1, \dots, T_{\mathrm{total}}$}
  \State $x_t \gets (1-\alpha)x_{t-1} + \alpha f(W_{\mathrm{res}}x_{t-1} + W_{\mathrm{in}}u_t + b)$
\EndFor
\State discard $t \le T_w$ (washout); stack remaining training states into $X$, targets into $y$
\State $W_{\mathrm{out}} \gets \left(X^{\!\top}X + \lambda\id\right)^{-1}X^{\!\top}y$ \Comment{$\lambda$ chosen on validation data}
\State \Return predictions $\hat y_t = W_{\mathrm{out}}\tilde x_t$ on the untouched test span
\end{algorithmic}
\end{algorithm*}

\section{What makes a good reservoir: the formal properties}
\label{sec:properties}

Three properties, introduced for ESNs but stated here in a form that transfers
verbatim to quantum reservoirs, separate systems that can compute from systems
that merely fluctuate.

\subsection{Echo state property}

\begin{definition}[Echo state property]
\label{def:esp}
A driven system with update $x_{t+1} = F(x_t, u_{t+1})$ on a compact state set
has the \emph{echo state property} (ESP) if for every input sequence
$\{u_t\}$ and every pair of initial conditions $x_0, x_0'$, the resulting
trajectories converge: $\lVert x_t - x_t'\rVert \to 0$ as $t \to \infty$.
Equivalently, after a transient the state is a function of the input history
alone, $x_t = \phi(\dots, u_{t-1}, u_t)$, independent of initialisation~\citep{jaeger2001echo}.
\end{definition}

The ESP is what makes the washout finite and the learned read-out
reproducible: without it, the ``features'' depend on an arbitrary initial
condition and the training problem is ill-posed. For the $\tanh$ ESN, a
sufficient condition is that the largest singular value of
$W_{\mathrm{res}}$ be below $1$; the everyday design rule of setting the
\emph{spectral radius} $\rho_{\mathrm{sr}}(W_{\mathrm{res}})$ slightly below
$1$ is a necessary condition (for zero input) and works well in practice,
though it is not sufficient in general~\citep{jaeger2001echo,lukosevicius2009reservoir}.
Extensions of the ESP to non-stationary systems and to subsystems --- needed when
only part of a larger system is observed, as is typical for quantum
reservoirs --- are developed in~\citet{kobayashi2024extending}.

\subsection{Fading memory}

\begin{definition}[Fading memory]
\label{def:fading}
A causal, time-invariant operator $M$ mapping input sequences to outputs has
\emph{fading memory} if inputs that agree on a recent window and differ only in
the remote past produce outputs that are close: for every $\varepsilon>0$
there exist $\delta>0$ and a decreasing weight sequence $w_k \downarrow 0$ such
that $\sup_k w_k\,|u_{t-k}-u'_{t-k}| < \delta$ implies
$|M[u]_t - M[u']_t| < \varepsilon$~\citep{boyd1985fading}.
\end{definition}

Fading memory formalises the intuition that a forecaster should care about the
recent past more than the distant past, with influence decaying continuously.
\citet{boyd1985fading} proved that operators with fading memory are exactly
those approximable by Volterra series --- the classical result that underlies all
reservoir universality theorems. In a reservoir with the ESP, fading memory is
automatic: the same contraction that erases the initial condition also erases
the remote input past. The tension at the heart of reservoir design is that
\emph{too much} contraction erases the recent past as well; the useful regime
lies near, but on the stable side of, the boundary
(Section~\ref{sec:edge-of-chaos}).

\subsection{Separation and the memory--nonlinearity trade-off}

The read-out is linear, so any two input histories that must be mapped to
different outputs must first be mapped, by the reservoir, to linearly
distinguishable states. This \emph{separation property}~\citep{maass2002realtime}
is the reservoir's expressivity requirement, complementary to the ESP's
stability requirement. A quantitative unification is the information-processing
capacity of \citet{dambre2012information}: for a reservoir with $N_x$
linearly independent state variables driven by i.i.d.\ input, the total
capacity to reconstruct \emph{any} orthonormal set of functions of the input
history (linear or nonlinear) is bounded by $N_x$, and saturates it for
fading-memory systems. Capacity spent on nonlinear functions is capacity not
available for linear memory --- the \emph{memory--nonlinearity trade-off} that
every reservoir, classical or quantum, must budget for the task at hand.

\subsection{Linear memory capacity}
\label{sec:memory-capacity}

The standard operational probe of memory is Jaeger's \emph{memory
capacity}~\citep{jaeger2001echo}. Drive the reservoir with i.i.d.\ scalar input
$u_t$, train a separate linear read-out to reconstruct each delayed input
$y_t = u_{t-d}$, and define
\begin{equation}
\begin{aligned}
  \mathrm{MC}_d \;&=\; \max_{w}\;
  \frac{\operatorname{cov}^2\!\left(u_{t-d},\, w^{\!\top}\tilde x_t\right)}
       {\operatorname{var}(u_{t-d})\,\operatorname{var}\!\left(w^{\!\top}\tilde x_t\right)},\\
  \mathrm{MC} \;&=\; \sum_{d=1}^{\infty} \mathrm{MC}_d .
\end{aligned}
  \label{eq:mc}
\end{equation}
Each $\mathrm{MC}_d \in [0,1]$ is the fraction of the delay-$d$ input variance
linearly recoverable from the state; $\mathrm{MC} \le N_x$ with equality
approached by linear reservoirs~\citep{jaeger2001echo,dambre2012information}.
In this document $\mathrm{MC}$ is estimated with delays up to $d_{\max}=25$ and
reported alongside every reservoir comparison (Part~VII), because a forecasting
model's usable history is bounded by it.

\subsection{Universality}
\label{sec:universality}

The formal justification for training only a linear read-out is a family of
universality theorems. For ESNs, \citet{grigoryeva2018echo} proved that echo
state networks with linear read-outs are universal: any causal, time-invariant
fading-memory filter on uniformly bounded inputs can be approximated to
arbitrary accuracy by an ESN of sufficient size. The proof route --- show the
family of reservoir functionals forms a polynomial algebra that separates
points, then apply Stone--Weierstrass --- is worth internalising because it
transfers: universality proofs for quantum reservoirs
(Part~II, Section~\ref{sec:qrc-universality}) follow the same architecture,
with the algebra generated by measured observables of the driven quantum
state~\citep{chen2019learning,nokkala2021gaussian,martinez2023finite}. Two
caveats temper the theorem's practical force. Universality is an existence
statement about \emph{some} sufficiently large reservoir, not a statement that
a random one of fixed size is good; and it says nothing about sample
efficiency. Both gaps are exactly where empirical reservoir quality --- and any
prospective quantum edge --- lives.

\section{Operating point: the edge of chaos}
\label{sec:edge-of-chaos}

For a fixed task, reservoir performance as a function of the dynamical regime
is strikingly non-monotonic. Deep in the ordered (strongly contracting) regime
the state barely responds to input and memory is short; deep in the chaotic
regime the ESP fails and the state amplifies microscopic differences until the
input signal is scrambled. Performance typically peaks near the
order--chaos boundary --- the \emph{edge of chaos} --- where correlations are long
but reproducibility survives~\citep{boedecker2012information}. For ESNs the
practical dials are the spectral radius $\rho_{\mathrm{sr}}$, the input
scaling, and the leak rate $\alpha$; systematic constraints on their joint
choice are mapped in~\citet{storm2022constraints}. Two further classical
lessons recur in the quantum setting. First, reservoirs whose dynamics have a
symmetry the task must break (for the $\tanh$ ESN, the odd symmetry
$f(-z)=-f(z)$) fail on symmetric tasks unless the symmetry is explicitly
broken, e.g.\ by biases~\citep{herteux2020breaking}; the quantum analogue ---
symmetries of the reservoir Hamiltonian polluting chaos diagnostics and
constraining which observables carry information --- appears in Part~VIII.
Second, the edge is a property of the \emph{driven} system, not the autonomous
one: input strength shifts it, so the operating point must be tuned with the
data in the loop. In quantum many-body reservoirs the edge acquires a second,
temporal dimension (the injection interval relative to the Thouless time),
identified by \citet{kobayashi2026edge} and reproduced numerically in Part~VII
of this document.

\section{Reservoir computing for chaotic and geophysical systems}
\label{sec:rc-chaos}

The application benchmark that made modern RC famous is model-free prediction
of chaos. \citet{jaeger2004harnessing} demonstrated orders-of-magnitude
improvements in Mackey--Glass~\citep{mackey1977oscillation} prediction;
\citet{pathak2018model} showed that a parallelised ESN can forecast the
Kuramoto--Sivashinsky equation for multiple Lyapunov times and reproduce its
climate (attractor statistics); hybrids of reservoirs with imperfect
knowledge-based models outperform either alone~\citep{pathak2018hybrid};
and \citet{arcomano2020machine} scaled the approach to a global atmospheric
model, establishing RC as a serious tool for the weather--climate problem class
this document targets. Careful comparisons find well-tuned reservoirs
competitive with, and far cheaper to train than, backpropagated recurrent
networks on these tasks~\citep{vlachas2020backpropagation}. Chaotic benchmarks
(Lorenz-63~\citep{lorenz1963deterministic}, Mackey--Glass) therefore serve in
this document as controlled stand-ins for the harder, noisier renewable and
climate series of Part~VI, and the NARMA family serves as the standard
memory-plus-nonlinearity stress test.

\section{Baseline discipline}
\label{sec:baselines}

Because Parts~VI--VIII evaluate quantum reservoirs against economically
meaningful forecasting tasks, the classical baselines are fixed here once and
used everywhere. (i)~\emph{Mean predictor}: $\NMSE = 1$ by construction; drawn
as a horizontal line on every error axis. (ii)~\emph{Persistence}:
$\hat y_{t+h} = y_t$; embarrassingly strong on slowly varying series such as
solar irradiance at short horizons, and the baseline most often missing from
overclaiming studies. (iii)~\emph{Linear auto-regression on $p$ lags}: the
cheapest model with memory; any nonlinear reservoir must beat it to justify
its existence. (iv)~\emph{Size-matched ESN}: Eqs.~\eqref{eq:leaky} and
\eqref{eq:ridge-solution} with the number of read-out features equal to the
quantum feature count, so comparisons are at matched read-out capacity.
A quantum reservoir result in this document is reported as interesting only
relative to all four, and Part~VIII adds a fifth, quantum-internal control:
a structureless (Haar-random) quantum reservoir of the same size.
